% 第 12 题三棱锥三视图：正视图（左上，等腰三角形）、侧视图（右上，直角三角形，直角在左下）、
% 俯视图（正视图正下方，倒置等腰三角形）。
% 坐标以正视图底边左端为原点、底边为 y=0、向上为正；单位 u = 41.09 pt = 1.4442 cm
% （原图把「1」这一长度画成 41.09 pt）；三视图形状：底 2√3、高 1。
% 尺寸线：竖直尺寸线在图形左侧 0.2u 处（两端外向箭头，线身被数字 1 断开，上下端各 ±0.1u 横杠）；
% 水平尺寸线在底边（俯视图为顶边）外侧 0.2u 处（两端外向箭头，被数字断开，边界处 ±0.1u 竖直端线）；
% 尺寸线按标准画法，箭头尖对齐尺寸界线。图注在图形正下方。
\begin{tikzpicture}[>=stealth,x=1.4442cm,y=1.4442cm,every node/.style={font=\small}]
\def\s{1.732}   % √3（半底宽）
% ===== 正视图（左上）=====
\draw (0,0)--(\s,1)--(2*\s,0)--cycle;
\draw (\s,1)--(\s,0);
% 左侧竖直尺寸线
\draw[->] (-0.268,0.30)--(-0.268,0);
\draw[->] (-0.268,0.70)--(-0.268,1);
\node at (-0.268,0.5) {\(1\)};
\draw (-0.368,1)--(-0.168,1);
\draw (-0.368,0)--(-0.168,0);
% 底边下方两条水平尺寸线（各标半底 √3）
\draw[->] ({0.5*\s-0.2},-0.2)--(0,-0.2);
\draw[->] ({0.5*\s+0.2},-0.2)--(\s,-0.2);
\node at ({0.5*\s},-0.2) {\(\sqrt{3}\)};
\draw[->] ({1.5*\s-0.2},-0.2)--(\s,-0.2);
\draw[->] ({1.5*\s+0.2},-0.2)--(2*\s,-0.2);
\node at ({1.5*\s},-0.2) {\(\sqrt{3}\)};
\draw (0,-0.1)--(0,-0.3);
\draw (\s,-0.1)--(\s,-0.3);
\draw (2*\s,-0.1)--(2*\s,-0.3);
\node[anchor=north] at (\s,-0.46) {正视图};
% ===== 侧视图（右上，与正视图同基线）=====
\draw (4.732,0)--(5.732,0)--(4.732,1)--cycle;
% 左侧竖直尺寸线
\draw[->] (4.532,0.30)--(4.532,0);
\draw[->] (4.532,0.70)--(4.532,1);
\node at (4.532,0.5) {\(1\)};
\draw (4.432,1)--(4.632,1);
\draw (4.432,0)--(4.632,0);
% 底边下方水平尺寸线（标 1）
\draw[->] (5.032,-0.2)--(4.732,-0.2);
\draw[->] (5.432,-0.2)--(5.732,-0.2);
\node at (5.232,-0.2) {\(1\)};
\draw (4.732,-0.1)--(4.732,-0.3);
\draw (5.732,-0.1)--(5.732,-0.3);
\node[anchor=north] at (5.232,-0.46) {侧视图};
% ===== 俯视图（正视图正下方，同宽同高）=====
\draw (0,-1.2)--(2*\s,-1.2)--(\s,-2.2)--cycle;
\draw (\s,-1.2)--(\s,-2.2);
% 左侧竖直尺寸线
\draw[->] (-0.268,-1.50)--(-0.268,-1.2);
\draw[->] (-0.268,-1.90)--(-0.268,-2.2);
\node at (-0.268,-1.7) {\(1\)};
\draw (-0.368,-1.2)--(-0.168,-1.2);
\draw (-0.368,-2.2)--(-0.168,-2.2);
% 顶边上方两条水平尺寸线（各标半底 √3，与正视图镜像）
\draw[->] ({0.5*\s-0.2},-1.0)--(0,-1.0);
\draw[->] ({0.5*\s+0.2},-1.0)--(\s,-1.0);
\node at ({0.5*\s},-1.0) {\(\sqrt{3}\)};
\draw[->] ({1.5*\s-0.2},-1.0)--(\s,-1.0);
\draw[->] ({1.5*\s+0.2},-1.0)--(2*\s,-1.0);
\node at ({1.5*\s},-1.0) {\(\sqrt{3}\)};
\draw (0,-0.9)--(0,-1.1);
\draw (\s,-0.9)--(\s,-1.1);
\draw (2*\s,-0.9)--(2*\s,-1.1);
\node[anchor=north] at (\s,-2.46) {俯视图};
\end{tikzpicture}
